\documentclass[11pt,letterpaper]{article}

\usepackage[spanish,es-noshorthands]{babel}
\usepackage{fontspec}
\setmainfont{Source Serif Pro}
\usepackage{microtype}
\usepackage{geometry}
\usepackage{booktabs}
\usepackage{tabularx}
\usepackage{enumitem}
\usepackage{xcolor}
\usepackage{graphicx}
\usepackage{csquotes}
\usepackage{amsmath}
\usepackage{siunitx}
\usepackage{tikz}
\usetikzlibrary{arrows.meta,positioning,shapes.geometric,fit,backgrounds,patterns}
\usepackage{xurl}
\usepackage[hidelinks]{hyperref}
\usepackage[backend=biber,style=apa,sorting=nyt]{biblatex}
\DeclareLanguageMapping{spanish}{spanish-apa}

\geometry{margin=2.3cm}
% Los rótulos de las figuras usan nodos de ancho fijo con texto alineado a la
% derecha, y la bibliografía compone en modo sloppy: ambos generan avisos de
% línea holgada que son cosméticos y no indican un problema de composición.
% Se silencian para que `make check` siga siendo sensible a los desbordes reales.
\hbadness=10000
\setlength{\emergencystretch}{1em}
\setlength{\parindent}{0pt}
\setlength{\parskip}{0.55em}
\setlist{nosep}
% Las figuras ocupan entre el 60 y el 70 por ciento de la caja de texto.
% Sin estos parámetros LaTeX las difiere todas al final del documento.
\renewcommand{\topfraction}{0.92}
\renewcommand{\bottomfraction}{0.85}
\renewcommand{\textfraction}{0.06}
\renewcommand{\floatpagefraction}{0.72}
\setcounter{topnumber}{2}
\setcounter{totalnumber}{3}
\addbibresource{referencias.bib}
\AtBeginBibliography{\sloppy}

\sisetup{
  output-decimal-marker = {,},
  group-separator = {.},
  group-minimum-digits = 4,
  detect-all
}

% Colores: Anthropic, OpenAI, tercer jugador (Google), producto de referencia
% y señal de alerta o dato dominado
\definecolor{cAnt}{RGB}{186,98,52}
\definecolor{cOAI}{RGB}{16,110,104}
\definecolor{cGoo}{RGB}{116,66,150}
\definecolor{cRef}{RGB}{70,90,112}
\definecolor{cAviso}{RGB}{176,42,42}

% Barra horizontal con rótulo a la izquierda y valor a la derecha.
% #1 y, #2 fin de la barra, #3 color, #4 rótulo, #5 valor, #6 x del valor
\newcommand{\barrarot}[6]{%
  \node[left, text width=4.2cm, align=right, font=\scriptsize] at (-0.15,#1) {#4};%
  \fill[#3] (0,#1-0.17) rectangle (#2,#1+0.17);%
  \node[right, font=\scriptsize] at (#6,#1) {#5};%
}

\title{\vspace{-1.2cm}\textbf{Anthropic y OpenAI: dos ofensivas tecnológicas
sobre el mismo mercado}\\
\large Clasificación competitiva y gráfico vectorial con datos del 1 de
septiembre de 2026}
\author{Carlos Luis Rojas Aragonés\\
\small Gestión del Desarrollo Tecnológico}
\date{\small 1 de septiembre de 2026}

\begin{document}

\maketitle
\vspace{-0.6em}

\section{El sector y por qué lo elegí}

El mercado que analizo es el de los \textbf{modelos fundacionales de lenguaje
vendidos como servicio}, es decir, la capacidad de razonamiento entregada por
interfaz de programación y facturada por millón de tokens, junto con las
aplicaciones que se construyen encima. Cumple con exactitud la condición del
enunciado: es un mercado cuya oferta, su precio y hasta su definición dependen
del estado de desarrollo del sistema tecnológico de unos pocos actores. No hay
un sustituto no tecnológico. Cuando Anthropic recorta el precio de la lectura de
caché un 75 por ciento o cuando OpenAI publica una generación con un 54 por
ciento más de eficiencia por token, el mercado entero se recoloca en cuestión de
semanas.

Los dos competidores que comparo son \textbf{Anthropic} y \textbf{OpenAI}. Los
elijo porque comparten origen (siete de los ocho fundadores de Anthropic salieron
de OpenAI en 2021), compiten con el mismo insumo, el mismo tipo de cliente y una
tecnología de base casi idéntica, y sin embargo han acabado con dos estrategias
distinguibles y con dos repartos de mercado casi opuestos. Es el caso de libro
que pide el foro: estrategias algo distintas sobre un mismo mercado.

Anticipo la tesis, porque estructura todo lo que sigue: \textbf{no se reparten un
mercado, se reparten dos}. Anthropic domina el gasto empresarial por interfaz de
programación (40 por ciento) y muy en particular la codificación (54 por ciento);
OpenAI domina el consumo (46,4 por ciento de la audiencia de asistentes) y ha
convertido esa audiencia en un negocio de publicidad y comercio que Anthropic ni
siquiera intenta \parencite{menlo2025,sensortower2026}. La comparación técnica
que hago en las secciones \ref{sec:benchmarks} a \ref{sec:equipo} explica por qué ese reparto no es casualidad:
está inscrito en la arquitectura de precios y en el perfil de capacidades de cada
uno.

Una advertencia de método antes de empezar. \textbf{No busco un ganador.} Con los
datos que manejo, ninguna de las dos compañías domina a la otra en el sentido
técnico del término, y precisamente eso es lo que hace interesante el gráfico
vectorial de la sección \ref{sec:vectorial}: la frontera de eficiencia tiene cuatro puntos y está
repartida dos a dos.

\section{Marco de análisis: qué categorías uso y qué mido}

\subsection*{Las categorías de clasificación}

El enunciado pide clasificar a las dos compañías \enquote{a partir de las
categorías anteriores}. Interpreto que se refiere a la taxonomía clásica de
estrategias tecnológicas de \textcite{freeman1997}, que distingue seis posturas
(ofensiva, defensiva, imitativa, dependiente, tradicional y oportunista), y la
cruzo con el posicionamiento por figuras de mérito propio de la hoja de ruta
tecnológica \parencite{deweck2022}, que es el instrumento que el gráfico
vectorial formaliza. Añado dos ejes complementarios que en este sector resultan
decisivos y que la literatura tiene bien establecidos: el de los \emph{activos
complementarios} de \textcite{teece1986}, que explica quién se apropia del valor
de una innovación, y el de las \emph{guerras de estándares} de
\textcite{shapiro1999}, que explica por qué una empresa regala un protocolo.
Si el módulo se refería a otro juego de categorías, la clasificación del
Cuadro~\ref{tab:clasificacion} sigue siendo válida en sus propios términos.

\subsection*{Las figuras de mérito}

Una figura de mérito es una medida cuantitativa del desempeño de una tecnología
en la función que presta, independiente de la implementación
\parencite{deweck2022}. En este mercado el cliente compra dos cosas a la vez, y
por eso hacen falta dos ejes:

\begin{enumerate}
  \item \textbf{Valor entregado}, medido como capacidad verificable en
    evaluaciones públicas e independientes. Uso el \emph{Artificial Analysis
    Intelligence Index}, un índice compuesto de veinticuatro evaluaciones que
    ejecuta un tercero sin vínculo con ninguno de los dos proveedores
    \parencite{aa2026modelos}, y lo contrasto con tres evaluaciones adicionales
    en la sección \ref{sec:benchmarks}.
  \item \textbf{Coste soportado por el cliente}, medido como coste en dólares de
    una tarea tipo definida en tokens, calculado a partir de las tarifas
    oficiales publicadas por cada compañía
    \parencite{anthropic2026precios,openai2026precios}. No uso el precio por
    millón de tokens como figura de mérito, y explico en la sección \ref{sec:precios} por qué
    ese número, tomado suelto, es engañoso.
\end{enumerate}

Las otras dos dimensiones que pide la comparación, innovación y equipo técnico,
las trato como \emph{determinantes} de esos dos ejes y no como ejes propios: la
innovación explica el desplazamiento vertical del gráfico vectorial y el equipo
explica la velocidad con la que ese desplazamiento se produce.

\section{Clasificación de las dos compañías}

\begin{figure}[tbp]
  \centering
  \begin{tikzpicture}
    % ---------------- panel A: trayectoria 2021-2026 ----------------
    \node[font=\scriptsize\itshape, gray!35!black, right] at (0,3.35)
      {Panel A. Origen común y separación de trayectorias, 2021 a 2026};
    \draw[gray!60!black] (0,0.95) -- (13.20,0.95);
    \foreach \xv/\lab in {0/2021, 2.2/2022, 4.4/2023, 6.6/2024, 8.8/2025,
                          11.0/2026, 13.2/2027} {
      \draw[gray!60!black] (\xv,0.95) -- (\xv,0.82);
      \node[below, font=\tiny, gray!55!black, fill=white, inner sep=1.2pt]
        at (\xv,0.82) {\lab};
    }
    % --- pista superior: Anthropic ---
    \foreach \xv/\prof/\lab in {0.10/1.55/{fundación por\\siete salidas\\de
                                  OpenAI},
                                4.77/0.50/{Claude 1},
                                6.97/1.55/{familia\\Claude 3},
                                8.43/0.50/{publica MCP},
                                9.53/1.55/{Claude Code\\disponible},
                                11.00/0.50/{dona MCP a la\\Agentic AI
                                  Foundation},
                                11.92/1.90/{Fable 5 y\\Mythos 5},
                                12.47/0.85/{Fable 5.1}} {
      \draw[cAnt, thick] (\xv,0.95) -- (\xv,0.95+\prof);
      \fill[cAnt] (\xv,0.95+\prof) circle (1.6pt);
      \node[above, font=\tiny, cAnt, align=center, fill=white,
        inner sep=1.5pt] at (\xv,1.02+\prof) {\lab};
    }
    % --- pista inferior: OpenAI ---
    \foreach \xv/\prof/\lab in {4.03/0.50/{ChatGPT},
                                4.77/1.45/{GPT-4},
                                8.07/0.50/{serie o\\de razonamiento},
                                10.30/1.45/{reestructuración\\en sociedad de
                                  beneficio público},
                                10.62/0.50/{navegador\\Atlas},
                                12.10/1.90/{GPT-5.6 y\\ChatGPT Work}} {
      \draw[cOAI, thick] (\xv,0.95) -- (\xv,0.95-\prof);
      \fill[cOAI] (\xv,0.95-\prof) circle (1.6pt);
      \node[below, font=\tiny, cOAI, align=center, fill=white,
        inner sep=1.5pt] at (\xv,0.88-\prof) {\lab};
    }
    \draw[cAnt, thick] (0.15,-2.15) -- (0.75,-2.15);
    \node[right, font=\tiny, cAnt] at (0.85,-2.15) {Anthropic};
    \draw[cOAI, thick] (4.10,-2.15) -- (4.70,-2.15);
    \node[right, font=\tiny, cOAI] at (4.80,-2.15) {OpenAI};

    % ---------------- panel B: cadencia de lanzamiento de 2026 -------------
    \begin{scope}[shift={(0,-5.30)}]
      \node[font=\scriptsize\itshape, gray!35!black, right] at (0,2.55)
        {Panel B. Cadencia de lanzamiento de modelos entre el 1 de enero y el
         1 de septiembre de 2026};
      \draw[gray!60!black] (0,0.85) -- (13.20,0.85);
      \foreach \xv/\lab in {0/ene, 1.625/feb, 3.25/mar, 4.875/abr, 6.5/may,
                            8.125/jun, 9.75/jul, 11.375/ago, 13.0/sep} {
        \draw[gray!60!black] (\xv,0.85) -- (\xv,0.74);
        \node[below, font=\tiny, gray!55!black] at (\xv,0.74) {\lab};
      }
      \foreach \xv/\lab in {1.84/{4.6}, 2.49/{S4.6}, 5.20/{Myt.},
                            5.69/{4.7}, 7.96/{4.8}, 8.56/{F5}, 9.70/{S5},
                            10.95/{O5}, 13.00/{F5.1}} {
        \draw[cAnt, thick] (\xv,0.85) -- (\xv,1.55);
        \fill[cAnt] (\xv,1.55) circle (1.7pt);
        \node[above, font=\tiny, cAnt, rotate=52, anchor=west]
          at (\xv,1.62) {\lab};
      }
      \foreach \xv/\lab in {1.84/{5.3C}, 3.46/{5.4}, 4.11/{mini}, 6.06/{5.5},
                            6.71/{Inst.}, 9.47/{5.6p}, 10.19/{5.6}} {
        \draw[cOAI, thick] (\xv,0.85) -- (\xv,0.15);
        \fill[cOAI] (\xv,0.15) circle (1.7pt);
        \node[below, font=\tiny, cOAI, rotate=52, anchor=east]
          at (\xv,0.08) {\lab};
      }
      \node[right, font=\tiny, cAnt, align=left] at (0.15,1.95)
        {9 lanzamientos en 8 meses};
      \node[right, font=\tiny, cOAI, align=left] at (0.15,-0.60)
        {7 lanzamientos en 8 meses};
    \end{scope}
  \end{tikzpicture}
  \caption{Origen común, trayectorias distintas y cadencia casi idéntica.
  Panel A: Anthropic se funda en enero de 2021 con siete antiguos empleados de
  OpenAI \parencite{wikipedia2026anthropic}; los hitos de la pista inferior son
  de OpenAI \parencite{wikipedia2026openai}. Panel B: fechas de lanzamiento de
  2026 según los registros públicos de cada familia de modelos
  \parencite{wikipedia2026claude,wikipedia2026gpt56}. Abreviaturas de la pista
  superior: 4.6 y 4.7 y 4.8 son versiones de Claude Opus, S4.6 y S5 de Claude
  Sonnet, Myt. es la vista previa de Mythos, F5 y F5.1 son Claude Fable y O5 es
  Claude Opus 5. En la inferior: 5.3C es GPT-5.3-Codex, mini es la pareja
  mini y nano de GPT-5.4, Inst. es GPT-5.5 Instant, 5.6p es la vista previa
  restringida de GPT-5.6 y 5.6 su disponibilidad general. Elaboración propia.}
  \label{fig:cronologia}
\end{figure}

El Panel B de la Figura~\ref{fig:cronologia} contiene el primer resultado
objetivo de este trabajo y conviene enunciarlo antes de clasificar nada:
\textbf{la cadencia de lanzamiento de las dos compañías es prácticamente la
misma}, nueve frente a siete eventos en ocho meses, con un intervalo medio de
veintisiete y treinta y cuatro días respectivamente. Quien espere encontrar
aquí un innovador rápido frente a un seguidor lento no va a encontrarlo. La
diferencia no está en el ritmo, está en la dirección.

El Cuadro~\ref{tab:clasificacion} sitúa a las dos compañías en las cuatro
taxonomías del marco.

\begin{table}[tbp]
  \centering
  \small
  \caption{Clasificación de los dos jugadores según cuatro juegos de
  categorías.}
  \label{tab:clasificacion}
  \begin{tabularx}{\textwidth}{@{}>{\raggedright\arraybackslash}p{0.17\textwidth}
    >{\raggedright\arraybackslash}X >{\raggedright\arraybackslash}X@{}}
    \toprule
    \textbf{Categoría} & \textbf{Anthropic} & \textbf{OpenAI} \\
    \midrule
    Estrategia tecnológica \parencite{freeman1997}
      & Ofensiva. Investigación propia, publicación científica, ambición de
        estar en la frontera y no detrás de ella
      & Ofensiva. Idéntica postura, con el añadido de haber creado la categoría
        de producto en 2022 \\
    Variante de la ofensiva
      & \emph{De profundidad}: gana en un dominio acotado (trabajo agente de
        larga duración y codificación) y se expande desde ahí
      & \emph{De amplitud}: cubre todas las modalidades y todos los públicos, y
        acepta perder profundidad en cada uno \\
    Apropiación del valor \parencite{teece1986}
      & Integración hacia atrás y lateral: proveedores de cómputo múltiples y
        distribución por las tres nubes ajenas
      & Integración hacia adelante: aplicación de consumo, navegador,
        publicidad, comercio y hardware propio \\
    Postura ante el estándar \parencite{shapiro1999}
      & Abierta. Publica el protocolo de contexto (MCP) en 2024 y lo dona a una
        fundación en enero de 2026
      & Adopción tardía. Sustituye su interfaz propietaria de agentes por MCP y
        entra como cofundador de la misma fundación \\
    Posición por figuras de mérito \parencite{deweck2022}
      & Líder en valor entregado por tarea; no líder en coste
      & Líder en coste por tarea en interacciones cortas; no líder en valor \\
    Modelo de ingresos
      & Alrededor del 80 por ciento empresa e interfaz de programación
      & Consumo mayoritario hasta agosto de 2026, cuando la empresa lo supera
        por primera vez \\
    \bottomrule
  \end{tabularx}
\end{table}

La conclusión de la clasificación es que \textbf{ambas son ofensivas, que es la
situación menos estable de la taxonomía de Freeman}, porque obliga a las dos a
sostener el gasto de investigación sin poder amortizarlo imitando a nadie. Lo
que las separa no es la categoría, es el vector de la ofensiva. Y ahí sí hay dos
elecciones deliberadas y opuestas: Anthropic apuesta por la profundidad en el
trabajo agente y por un estándar abierto que le da distribución sin poseer el
canal; OpenAI apuesta por la amplitud y por poseer el canal entero, del navegador
al dispositivo.

\section{Primera dimensión: los \emph{benchmarks} públicos}
\label{sec:benchmarks}

\begin{figure}[tbp]
  \centering
  \begin{tikzpicture}
    % ---------------- panel A: índice compuesto independiente -------------
    \node[font=\scriptsize\itshape, gray!35!black, right] at (-4.35,3.30)
      {Panel A. Artificial Analysis Intelligence Index, esfuerzo máximo, misma
       versión del índice};
    \foreach \xv/\lab in {0/0, 1.8/20, 3.6/40, 5.4/60, 6.3/70} {
      \draw[gray!22] (\xv,-0.30) -- (\xv,2.90);
      \node[below, font=\tiny, gray!55!black] at (\xv,-0.30) {\lab};
    }
    \barrarot{2.60}{5.94}{cAnt}{Claude Fable 5.1}{66}{6.04}
    \barrarot{2.08}{5.67}{cAnt!80}{Claude Opus 5}{63}{6.04}
    \barrarot{1.56}{5.58}{cAnt!60}{Claude Fable 5}{62}{6.04}
    \barrarot{1.04}{5.49}{cOAI}{GPT-5.6 Sol}{61}{6.04}
    \barrarot{0.52}{5.13}{cRef!70}{Claude Opus 4.8 (referencia)}{57}{6.04}
    \barrarot{0.00}{5.04}{cRef!45}{GPT-5.5 (referencia)}{56}{6.04}

    % ---------------- panel B: ARC-AGI-3 ----------------
    \begin{scope}[shift={(0,-4.35)}]
      \node[font=\scriptsize\itshape, gray!35!black, right] at (-4.35,2.75)
        {Panel B. ARC-AGI-3, razonamiento interactivo sobre entornos nuevos, en
         por ciento};
      \foreach \xv/\lab in {0/0, 1.8/10, 3.6/20, 5.4/30, 6.3/35} {
        \draw[gray!22] (\xv,-0.30) -- (\xv,2.35);
        \node[below, font=\tiny, gray!55!black] at (\xv,-0.30) {\lab};
      }
      \barrarot{2.08}{5.44}{cAnt!80}{Claude Opus 5}{30,2\,\%}{6.04}
      \barrarot{1.56}{1.40}{cOAI}{GPT-5.6 Sol}{7,8\,\%}{6.04}
      \barrarot{1.04}{0.27}{cRef!70}{Claude Opus 4.8}{1,5\,\%}{6.04}
      \barrarot{0.52}{0.14}{cOAI!55}{GPT-5.6 Terra}{0,8\,\%}{6.04}
      \barrarot{0.00}{0.07}{cRef!45}{GPT-5.5}{0,4\,\%}{6.04}
    \end{scope}

    % ---------------- panel C: Terminal-Bench 2.1 ----------------
    \begin{scope}[shift={(0,-8.55)}]
      \node[font=\scriptsize\itshape, gray!35!black, right] at (-4.35,2.75)
        {Panel C. Terminal-Bench 2.1, tareas de agente en línea de órdenes,
         escala ampliada de 86 a 93 por ciento};
      \foreach \xv/\lab in {0/86, 1.575/88, 3.15/90, 4.725/92, 5.5/93} {
        \draw[gray!22] (\xv,-0.30) -- (\xv,2.35);
        \node[below, font=\tiny, gray!55!black] at (\xv,-0.30) {\lab};
      }
      \barrarot{2.08}{4.65}{cOAI!75}{GPT-5.6 Sol Ultra}{91,9\,\%}{6.04}
      \barrarot{1.56}{2.44}{cAnt!80}{Claude Opus 5}{89,1\,\%}{6.04}
      \barrarot{1.04}{2.20}{cOAI}{GPT-5.6 Sol}{88,8\,\%}{6.04}
      \barrarot{0.52}{1.58}{cAnt!55}{Claude Mythos 5}{88,0\,\%}{6.04}
      \barrarot{0.00}{1.58}{cRef!45}{GPT-5.5}{88,0\,\%}{6.04}
    \end{scope}
  \end{tikzpicture}
  \caption{Tres evaluaciones públicas con tres resultados distintos. Panel A:
  índice compuesto de veinticuatro evaluaciones medido por un tercero
  independiente; los valores proceden de la misma versión del índice, que es la
  única forma de compararlos entre sí
  \parencite{aa2026modelos,aa2026comparacion}. Panel B: puntuaciones publicadas
  de ARC-AGI-3, la evaluación de razonamiento sobre entornos que el modelo no ha
  visto nunca \parencite{arcagi3,benchlm2026arc}. Panel C: Terminal-Bench 2.1
  según los datos de lanzamiento de cada proveedor
  \parencite{openai2026sol,anthropic2026opus5}; la escala está ampliada porque
  los cinco modelos caben en siete puntos porcentuales, que es en sí mismo el
  resultado. Elaboración propia.}
  \label{fig:benchmarks}
\end{figure}

La Figura~\ref{fig:benchmarks} es deliberadamente incómoda: las tres
evaluaciones dan tres respuestas diferentes, y esa es la primera conclusión.

\textbf{En el índice compuesto independiente, Anthropic va delante por poco.}
Claude Fable 5.1 marca 66 puntos y GPT-5.6 Sol marca 61, con Claude Opus 5 en 63
\parencite{aa2026modelos,aa2026comparacion}. La distancia es de cinco puntos
sobre un índice cuyo recorrido útil en los últimos doce meses ha sido de unos
diez, de modo que equivale a poco más de medio año de progreso. Conviene además
recordar que el orden se ha invertido varias veces en 2026: en el análisis del 9
de julio, GPT-5.6 Sol quedaba a un punto de Claude Fable 5
\parencite{aa2026gpt56}.

\textbf{En razonamiento sobre entornos nuevos, la distancia es de otro orden.}
En ARC-AGI-3, la evaluación diseñada expresamente para que no se pueda memorizar,
Claude Opus 5 obtiene un 30,2 por ciento y GPT-5.6 Sol un 7,8 por ciento, casi
cuatro veces menos; el resto de la industria está por debajo del 1,5 por ciento
\parencite{benchlm2026arc}. Es la única de las tres evaluaciones donde hay una
diferencia cualitativa y no de matiz, y merece leerse con cuidado: en marzo de
2026 todos los modelos de frontera puntuaban por debajo del 0,4 por ciento, de
modo que lo que mide el panel es un salto muy reciente y muy concentrado en un
solo modelo.

\textbf{En trabajo de agente sobre terminal, empatan y gana OpenAI en la
configuración de esfuerzo máximo.} Los cinco modelos del Panel C caben en siete
puntos porcentuales, con GPT-5.6 Sol Ultra en cabeza con un 91,9 por ciento
frente al 89,1 de Claude Opus 5 \parencite{openai2026sol,anthropic2026opus5}. Y
en el índice de agentes de codificación de Artificial Analysis, versión 1.1,
GPT-5.6 Sol obtiene 80 puntos frente a los 77,2 de Claude Fable 5, con menos de
la mitad de tokens de salida y menos de la mitad de tiempo
\parencite{wikipedia2026gpt56,aa2026gpt56}. Es decir: en la evaluación de
codificación con criterio independiente, gana OpenAI, y gana además por
eficiencia.

Hay una cuarta observación, metodológica, que considero la más importante de esta
sección. \textbf{Las evaluaciones se renuevan más deprisa que los modelos, y eso
rompe la comparabilidad.} Anthropic reporta Terminal-Bench 4.0 para Fable 5.1
(55,8 por ciento) mientras OpenAI reporta Terminal-Bench 2.1 para Sol (88,8 por
ciento) \parencite{anthropic2026fable51,openai2026sol}: son escalas distintas y
ponerlas en la misma barra sería un error grosero. El propio índice de
Artificial Analysis cambia de versión: Claude Opus 4.8 se anunció con 61,4 puntos
en mayo \parencite{aa2026opus48} y en la versión vigente del índice figura con 57
\parencite{aa2026comparacion}. Por eso todas las comparaciones de esta sección
son intraversión, y por eso el gráfico vectorial de la sección \ref{sec:vectorial} se construye
sobre una sola fuente y una sola versión del índice.

\section{Segunda dimensión: innovación}

Medir innovación con un solo número es imposible, así que uso cuatro
indicadores observables y los contrasto.

\textbf{Cadencia.} Ya está en el Panel B de la Figura~\ref{fig:cronologia}: nueve
lanzamientos frente a siete en ocho meses. Empate técnico.

\textbf{Creación de estándar.} Aquí la asimetría es grande y verificable.
Anthropic publicó en noviembre de 2024 el Model Context Protocol, la interfaz por
la que un modelo se conecta a herramientas y datos, y en enero de 2026 lo donó a
la Agentic AI Foundation bajo la Linux Foundation, con OpenAI y Block como
cofundadores y con AWS, Google, Microsoft, Cloudflare, GitHub y Bloomberg como
miembros \parencite{anthropic2026mcp}. OpenAI llegó a retirar su interfaz
propietaria de asistentes en favor del protocolo ajeno
\parencite{workos2026mcp}. Es un movimiento de manual de \textcite{shapiro1999}:
quien no controla el canal de distribución regala el estándar para que el canal
se construya alrededor de su forma de trabajar. Anthropic renunció al control del
protocolo y a cambio consiguió que el ecosistema de agentes de toda la industria
hable su idioma.

\textbf{Amplitud del frente de producto.} Aquí la asimetría es la contraria y de
magnitud parecida. Durante 2026 OpenAI ha abierto o consolidado la publicidad
dentro de ChatGPT (desde febrero, con un ritmo anualizado cercano a los mil
millones de dólares en agosto), el comercio dentro de la conversación con una
docena de grandes minoristas integrados, la fusión del navegador Atlas y del
agente de codificación en una sola aplicación de escritorio, y un dispositivo
físico previsto para 2027 \parencite{tnw2026friar,digitalapplied2026work}.
Anthropic no ha entrado en ninguno de esos cuatro frentes. Su ampliación de
producto ha sido lateral y siempre hacia el mismo cliente: Claude Code, Claude
Cowork, Claude Design y Claude Science \parencite{wikipedia2026claude}.

\textbf{Gobierno de la capacidad.} Este es el indicador que menos se cita y el
que más me interesa como CTO, porque las dos compañías han convergido en él por
caminos distintos. Anthropic separa desde abril de 2026 una línea de acceso
restringido, Mythos, para organizaciones verificadas de ciberseguridad y ciencias
de la vida, con el mismo modelo base que la línea general Fable pero con otro
régimen de salvaguardas \parencite{anthropic2026fable51,anthropic2026precios}.
OpenAI lanzó GPT-5.6 Sol el 26 de junio de 2026 en vista previa restringida a
socios cuya participación se comunicó al gobierno de los Estados Unidos, y
declaró que ese tipo de restricción no debería ser la norma
\parencite{techcrunch2026restriccion}. Trece días después el modelo estaba en
disponibilidad general \parencite{openai2026gpt56}.

Ese mismo indicador produjo en julio de 2026 el suceso más instructivo del año
para cualquiera que dirija un equipo técnico. OpenAI comunicó que dos de sus
modelos, entre ellos GPT-5.6 Sol, se salieron del entorno aislado de una
evaluación de capacidades ofensivas, alcanzaron infraestructura conectada a
internet y comprometieron sistemas de producción de un tercero para obtener las
soluciones de la propia evaluación \parencite{hackernews2026exploit}. La empresa
lo hizo público por iniciativa propia, lo que hay que reconocerle. Pero el hecho
tiene dos lecturas que afectan directamente a este trabajo: primera, que la
integridad de los \emph{benchmarks} de capacidad ofensiva ya no puede darse por
supuesta; segunda, que el aislamiento del entorno de evaluación es una figura de
mérito de ingeniería tan legítima como la puntuación que ese entorno produce.

\section{Tercera dimensión: precios}
\label{sec:precios}

\begin{figure}[tbp]
  \centering
  \begin{tikzpicture}
    % ---------------- panel A: tarifas de catálogo ----------------
    \node[font=\scriptsize\itshape, gray!35!black, right] at (-4.35,3.30)
      {Panel A. Tarifa oficial de salida, en dólares por millón de tokens,
       con la de entrada superpuesta};
    \foreach \xv/\lab in {0/0, 1.2/10, 2.4/20, 3.6/30, 4.8/40, 6.0/50} {
      \draw[gray!22] (\xv,-0.30) -- (\xv,2.95);
      \node[below, font=\tiny, gray!55!black] at (\xv,-0.30) {\lab};
    }
    \barrarot{2.60}{6.00}{cAnt}{Claude Fable 5.1}{50 y 10}{6.10}
    \barrarot{2.08}{3.00}{cAnt!80}{Claude Opus 5}{25 y 5}{6.10}
    \barrarot{1.56}{1.20}{cAnt!55}{Claude Sonnet 5}{10 y 2}{6.10}
    \barrarot{1.04}{2.40}{cOAI}{GPT-5.6 Sol}{20 y 4}{6.10}
    \barrarot{0.52}{1.44}{cOAI!70}{GPT-5.6 Terra}{12 y 2}{6.10}
    \barrarot{0.00}{0.14}{cOAI!45}{GPT-5.6 Luna}{1,2 y 0,2}{6.10}
    \foreach \y/\xe in {2.60/1.20, 2.08/0.60, 1.56/0.24, 1.04/0.48,
                        0.52/0.24, 0.00/0.024} {
      \fill[white, opacity=0.55] (0,\y-0.17) rectangle (\xe,\y+0.17);
      \draw[gray!45!black, thin] (\xe,\y-0.17) -- (\xe,\y+0.17);
    }

    % ---------------- panel B: coste de una tarea tipo ----------------
    \begin{scope}[shift={(0,-4.90)}]
      \node[font=\scriptsize\itshape, gray!35!black, right] at (-4.35,3.05)
        {Panel B. Coste de una tarea tipo, en dólares, con dos perfiles de uso};
      \foreach \xv/\lab in {0/0, 1.5/1, 3.0/2, 4.5/3, 6.0/4} {
        \draw[gray!22] (\xv,-0.35) -- (\xv,2.75);
        \node[below, font=\tiny, gray!55!black] at (\xv,-0.35) {\lab};
      }
      \foreach \y/\lab/\xl/\vl/\xc/\vc/\cs/\cc in {%
        2.40/{Claude Fable 5.1}/3.795/{2,53}/2.160/{1,44}/{cAnt}/{cAnt!40},
        1.80/{Claude Opus 5}/2.415/{1,61}/1.170/{0,78}/{cAnt!80}/{cAnt!32},
        1.20/{GPT-5.6 Sol}/3.414/{2,28}/0.936/{0,62}/{cOAI}/{cOAI!40},
        0.60/{GPT-5.6 Terra}/1.842/{1,23}/0.528/{0,35}/{cOAI!70}/{cOAI!28},
        0.00/{Claude Sonnet 5}/0.966/{0,64}/0.468/{0,31}/{cRef!60}/{cRef!25}} {
        \node[left, text width=4.2cm, align=right, font=\scriptsize]
          at (-0.15,\y) {\lab};
        \fill[\cs] (0,\y+0.03) rectangle (\xl,\y+0.23);
        \fill[\cc] (0,\y-0.23) rectangle (\xc,\y-0.03);
        \node[right, font=\tiny] at (\xl+0.08,\y+0.13) {\vl};
        \node[right, font=\tiny, gray!35!black] at (\xc+0.08,\y-0.13) {\vc};
      }
      \fill[gray!55] (0.15,-1.05) rectangle (0.70,-0.85);
      \node[right, font=\tiny, gray!25!black, align=left] at (0.78,-0.95)
        {perfil de agente: 1 M de tokens de entrada, 92\,\% en caché, 30 k de
         salida};
      \fill[gray!30] (0.15,-1.55) rectangle (0.70,-1.35);
      \node[right, font=\tiny, gray!25!black, align=left] at (0.78,-1.45)
        {perfil de conversación: 200 k de entrada, 80\,\% en caché, 20 k de
         salida};
    \end{scope}
  \end{tikzpicture}
  \caption{Dos arquitecturas de precio, no dos precios. Panel A: tarifas
  oficiales vigentes el 1 de septiembre de 2026; la barra sólida es la salida y
  el segmento claro del arranque es la entrada
  \parencite{anthropic2026precios,openai2026precios}. Las de GPT-5.6 son las
  promocionales de contexto corto, en vigor al menos hasta el 21 de noviembre de
  2026. Panel B: coste calculado por mí a partir de esas tarifas para dos
  perfiles de uso definidos en el texto; para OpenAI el perfil de agente entra en
  el tramo de contexto largo, que duplica la tarifa de entrada y sube la de
  salida. Elaboración propia.}
  \label{fig:precios}
\end{figure}

Comparar precios por millón de tokens es la trampa más común de este mercado, y
la Figura~\ref{fig:precios} está construida para desarmarla. Hay tres razones por
las que el número de catálogo engaña.

\textbf{Primera: los dos cobran el contexto largo de forma opuesta.} Anthropic
incluye la ventana completa de un millón de tokens a tarifa estándar, de modo que
una petición de 900.000 tokens se cobra al mismo precio unitario que una de 9.000
\parencite{anthropic2026precios}. OpenAI aplica un tramo de contexto largo por
encima de los 272.000 tokens que duplica la tarifa de entrada y sube la de salida
un 50 por ciento \parencite{openai2026precios}. Para una sesión de agente larga,
que es exactamente el caso de uso en disputa, esa diferencia estructural pesa más
que la tarifa nominal.

\textbf{Segunda: la caché es donde se está librando la competencia real.} El 1 de
septiembre de 2026 Anthropic bajó el precio de la lectura de caché de Fable de un
dólar a veinticinco centavos por millón, es decir al 2,5 por ciento de la tarifa
de entrada, frente al 10 por ciento habitual del resto de su catálogo y de la
competencia \parencite{anthropic2026precios,venturebeat2026fable51}. En el perfil
de agente del Panel B eso rebaja el coste de la tarea un 21 por ciento sin tocar
ni una sola tarifa de catálogo. Un cliente que mirase solo la columna de
\enquote{dólares por millón de tokens de salida} no habría visto nada.

\textbf{Tercera: el ranking de precio cambia de signo según el perfil de uso.}
Este es el resultado más útil de toda la sección. En el perfil de conversación,
GPT-5.6 Sol cuesta 0,62 dólares por tarea y Claude Opus 5 cuesta 0,78, un 26 por
ciento más. En el perfil de agente, Claude Opus 5 cuesta 1,61 dólares y GPT-5.6
Sol cuesta 2,28, un 41 por ciento más. \textbf{No hay un proveedor barato y otro
caro: hay dos estructuras de precio optimizadas para dos usos distintos}, y cada
una es la más barata en el suyo. El reparto del mercado de la sección \ref{sec:mercado} es, en
buena medida, la consecuencia contable de este hecho.

Queda una advertencia que afecta a todo lo anterior. Anthropic documenta que sus
modelos 4.7 y posteriores usan un tokenizador nuevo que produce aproximadamente un
30 por ciento más de tokens para el mismo texto \parencite{anthropic2026precios}.
Eso invalida la comparación de tarifas entre generaciones de Claude, y deja la
comparación entre Claude y GPT con un margen de error que no puedo cuantificar
porque no dispongo de una medición independiente de los dos tokenizadores sobre
el mismo corpus. Lo digo abiertamente: todos los costes del Panel B están
calculados en tokens, no en trabajo, y el trabajo es lo que el cliente compra.

\section{Cuarta dimensión: el equipo técnico}
\label{sec:equipo}

Sobre el equipo hay menos datos duros y más ruido, así que me limito a lo
verificable y señalo la calidad de cada fuente.

\textbf{Origen y liderazgo científico.} Anthropic la fundaron en enero de 2021
ocho personas procedentes de OpenAI, entre ellas Dario Amodei, que era
vicepresidente de investigación, su hermana Daniela, Jared Kaplan, coautor de las
leyes de escala \parencite{kaplan2020}, Chris Olah, que venía de fundar la línea
de interpretabilidad, Sam McCandlish, Tom Brown, primer autor del artículo de
GPT-3, Ben Mann y Jack Clark \parencite{wikipedia2026anthropic}. Es una escisión
inusual: no se llevaron un producto, se llevaron el núcleo de la agenda de
investigación. OpenAI, por su parte, conserva a Greg Brockman como presidente y
tiene como director científico a Jakub Pachocki desde mayo de 2024
\parencite{wikipedia2026openai}.

\textbf{Retención y flujo de talento.} Según el informe de talento de SignalFire,
la retención a dos años era del 80 por ciento en Anthropic, del 78 en Google
DeepMind y del 67 en OpenAI, y un ingeniero de OpenAI tenía ocho veces más
probabilidad de marcharse a Anthropic que a la inversa \parencite{signalfire2025}.
Son datos de 2025 y la edición de 2026 del mismo informe ya no los desglosa por
compañía \parencite{signalfire2026}, de modo que los trato como una dirección
razonablemente sólida y no como una foto actual. El episodio más citado en esa
dirección es la incorporación de Andrej Karpathy, cofundador de OpenAI, al equipo
de preentrenamiento de Anthropic el 19 de mayo de 2026
\parencite{cnbc2026karpathy,techcrunch2026karpathy}.

\textbf{Tamaño y productividad.} Anthropic declaraba del orden de 2.500 empleados
en 2026 y OpenAI unos 4.500 \parencite{wikipedia2026anthropic,wikipedia2026openai}.
Cruzando esas cifras con los ingresos anualizados de julio y agosto de 2026, 65.000
millones de dólares para Anthropic y 40.000 millones para OpenAI
\parencite{techcrunch2026ingresos,bloomberg2026openai}, salen 26,0 y 8,9 millones
de dólares de ingreso anualizado por empleado. Es una razón de 2,9 a 1. El cálculo
es mío y hay que tomarlo con pinzas, porque las dos compañías miden sus ingresos
de forma distinta y porque los recuentos de plantilla proceden de fuentes
terciarias, pero incluso admitiendo un error del cincuenta por ciento la
diferencia sigue siendo estructural y no de matiz.

\textbf{Y la contrapartida, que también hay que decirla.} La tensión entre
seguridad y comercialización no la resuelve ninguna de las dos. Jan Leike, que
había dejado OpenAI en 2024 por desacuerdos sobre seguridad y se había incorporado
a Anthropic, dimitió de Anthropic en febrero de 2026 con una advertencia pública
sobre la trayectoria del sector. La política de escalado responsable de Anthropic
va por su versión 3.1, de abril de 2026 \parencite{anthropic2026rsp}. Un equipo
excelente no es un equipo sin conflicto; es un equipo cuyo conflicto se hace
público.

\section{El gráfico vectorial}
\label{sec:vectorial}

\begin{table}[tbp]
  \centering
  \small
  \setlength{\tabcolsep}{4pt}
  \caption{Construcción de las dos coordenadas y de los vectores tecnológicos.
  El coste es el de la tarea de agente del Panel B de la
  Figura~\ref{fig:precios}. El producto de referencia es la media de los dos
  modelos de frontera vigentes el 28 de mayo de 2026: 56,5 puntos de índice y
  2,34 dólares por tarea.}
  \label{tab:vectores}
  \begin{tabularx}{\textwidth}{@{}l >{\raggedright\arraybackslash}X
    S[table-format=2.0] S[table-format=1.2]
    S[table-format=+2.1] S[table-format=+3.1]@{}}
    \toprule
    & \textbf{Vector tecnológico o comercial} & {\textbf{Índice}}
      & {\textbf{\$/tarea}} & {\textbf{$\Delta$ valor}}
      & {\textbf{$\Delta$ coste}} \\
    \midrule
    A1 & Fable 5: gama de cómputo de frontera (9 jun.) & 62 & 3.22
       & +5.5 & +37.6 \\
    A2 & Fable 5.1: lectura de caché al 2,5\,\% (1 sep.) & 66 & 2.53
       & +9.5 & +8.1 \\
    A3 & Opus 5: razonamiento adaptativo a mitad de tarifa (24 jul.)
       & 63 & 1.61 & +6.5 & -31.2 \\
    \addlinespace
    O1 & GPT-5.6 Sol al precio de la generación anterior (9 jul.)
       & 61 & 3.07 & +4.5 & +31.2 \\
    O2 & Recorte promocional de más del 20\,\% (21 ago.) & 61 & 2.28
       & +4.5 & -2.7 \\
    O3 & Gama Terra: equilibrio entre inteligencia y coste & 57 & 1.23
       & +0.5 & -47.5 \\
    O4 & Gama Luna: coste mínimo & 53 & 0.12 & -3.5 & -94.8 \\
    \bottomrule
  \end{tabularx}
\end{table}

\begin{figure}[tbp]
  \centering
  \begin{tikzpicture}[x=0.086cm, y=0.335cm,
      vec/.style  = {line width=1.05pt, -{Stealth[length=5.5pt]}},
      pto/.style  = {circle, fill, inner sep=2.1pt},
      cod/.style  = {font=\scriptsize\bfseries, inner sep=1.6pt},
      jugl/.style = {font=\footnotesize\bfseries, align=center, inner sep=2pt}]

    % rejilla
    \foreach \x in {-100,-80,-60,-40,-20,20,40}
      \draw[gray!20, thin] (\x,-5.6) -- (\x,11.6);
    \foreach \y in {-4,-2,2,4,6,8,10}
      \draw[gray!20, thin] (-104,\y) -- (46,\y);

    % frontera de eficiencia: banda ancha por detrás de todo lo demás
    \draw[gray!30, line width=3.2pt, line cap=round]
      (-94.8,-3.5) -- (-47.5,0.5) -- (-31.2,6.5) -- (8.1,9.5);
    \node[font=\tiny\itshape, text=gray!45!black, anchor=west, fill=white,
      inner sep=1.5pt] at (-88,2.2) {frontera de eficiencia};
    \draw[gray!45, thin] (-83,1.9) -- (-76,-0.9);

    % ejes
    \draw[thick,-{Stealth[length=6pt]}] (-104,0) -- (50,0);
    \draw[thick,-{Stealth[length=6pt]}] (0,-6.0) -- (0,12.4);
    \node[font=\footnotesize\bfseries, align=center, anchor=south]
      at (7,12.5) {$\Delta$ Valor entregado\\(puntos de índice)};
    \node[font=\footnotesize\bfseries, anchor=north east, fill=white,
      inner sep=2pt] at (46,-1.1)
      {$\Delta$ Coste por tarea de agente (por ciento)};

    % marcas de escala
    \foreach \x in {-100,-80,-60,-40,-20,20,40}
      \node[font=\tiny, text=gray!55!black, below=1pt] at (\x,0) {\x};
    \foreach \y in {-4,-2,2,4,6,8,10}
      \node[font=\tiny, text=gray!55!black, left=1pt] at (0,\y) {\y};

    % ---------- Anthropic ----------
    \draw[vec, cAnt] (0,0) -- (37.6,5.5)
      node[cod, text=cAnt, pos=0.55, above left=-2.5pt] {A1};
    \draw[vec, cAnt] (37.6,5.5) -- (8.1,9.5)
      node[cod, text=cAnt, pos=0.42, below left=-2.5pt] {A2};
    \draw[vec, cAnt] (0,0) -- (-31.2,6.5)
      node[cod, text=cAnt, pos=0.34, above right=-2.5pt] {A3};
    \node[pto, cAnt] at (37.6,5.5) {};
    \node[pto, cAnt] at (8.1,9.5) {};
    \node[pto, cAnt] at (-31.2,6.5) {};
    \node[jugl, text=cAnt, anchor=south] at (10,10.0) {Claude Fable 5.1};
    \node[font=\scriptsize, text=cAnt, anchor=west] at (39.8,6.1)
      {Fable 5};
    \node[jugl, text=cAnt, anchor=east] at (-34,6.9) {Claude\\Opus 5};

    % ---------- OpenAI ----------
    \draw[vec, cOAI] (0,0) -- (31.2,4.5)
      node[cod, text=cOAI, pos=0.80, below right=-2.5pt] {O1};
    \draw[vec, cOAI] (31.2,4.5) -- (-2.7,4.5)
      node[cod, text=cOAI, pos=0.62, above=-1pt] {O2};
    \draw[vec, cOAI] (-2.7,4.5) -- (-47.5,0.5)
      node[cod, text=cOAI, pos=0.55, below left=-2pt] {O3};
    \draw[vec, cOAI] (-47.5,0.5) -- (-94.8,-3.5)
      node[cod, text=cOAI, pos=0.45, below left=-2pt] {O4};
    \node[pto, cOAI] at (31.2,4.5) {};
    \node[pto, cOAI] at (-2.7,4.5) {};
    \node[pto, cOAI] at (-47.5,0.5) {};
    \node[pto, cOAI] at (-94.8,-3.5) {};
    \node[font=\scriptsize, text=cOAI, anchor=west] at (33.4,3.7)
      {Sol, precio inicial};
    \node[jugl, text=cOAI, anchor=east] at (-7.0,5.4) {GPT-5.6 Sol};
    \node[font=\scriptsize, text=cOAI, anchor=north] at (-47.5,-0.2)
      {Terra};
    \node[font=\scriptsize, text=cOAI, anchor=north] at (-94.8,-3.9)
      {Luna};

    % ---------- puntos dominados ----------
    \foreach \px/\py in {37.6/5.5, -2.7/4.5}
      \node[circle, draw=cAviso, thick, inner sep=3.6pt] at (\px,\py) {};
    \node[font=\tiny, text=cAviso, align=left, anchor=north west,
      fill=white, inner sep=2pt] at (-103,10.4) {Los dos puntos con anillo rojo están dominados:\\[-0.25em]
      Fable 5.1 supera a Fable 5 en valor y en coste, y\\[-0.25em]
      Claude Opus 5 supera a GPT-5.6 Sol en este perfil};

    % ---------- origen ----------
    \node[circle, draw, very thick, fill=white, inner sep=1.7pt] at (0,0) {};
    \node[font=\scriptsize, align=left, anchor=north east, fill=white,
      inner sep=2pt] at (46,-2.8)
      {Producto de referencia: la frontera del 28 de mayo de 2026\\[-0.25em]
       (Claude Opus 4.8 y GPT-5.5), 56,5 puntos y 2,34 \$ por tarea};
    \node[font=\scriptsize\itshape, text=gray!60!black, anchor=north west]
      at (-103,12.2) {Cuadrante favorable: arriba y a la izquierda};
  \end{tikzpicture}
  \caption{Gráfico vectorial de posicionamiento competitivo el 1 de septiembre
  de 2026. El origen es el producto de referencia del mercado tres meses atrás.
  Cada flecha es un vector del Cuadro~\ref{tab:vectores}: una decisión
  tecnológica o de tarifa con fecha conocida, no un desplazamiento postulado. Los dos
  puntos con anillo rojo están dominados por otro punto del gráfico. El
  eje vertical es el índice compuesto independiente del Panel A de la
  Figura~\ref{fig:benchmarks} \parencite{aa2026modelos,aa2026comparacion}; el
  horizontal, el coste calculado de la tarea de agente del Panel B de la
  Figura~\ref{fig:precios}. Las puntuaciones de índice de Terra y Luna están
  ajustadas por mí a la versión vigente sumando dos puntos a las publicadas el 9
  de julio \parencite{aa2026gpt56}. Elaboración propia.}
  \label{fig:vectorial}
\end{figure}

La Figura~\ref{fig:vectorial} es la respuesta gráfica que pide el enunciado, y
dice cinco cosas.

\textbf{Primera: la frontera de eficiencia tiene cuatro puntos y está repartida
dos a dos.} Claude Fable 5.1, Claude Opus 5, GPT-5.6 Terra y GPT-5.6 Luna no se
dominan entre sí: cada uno es el mejor en su tramo de disposición a pagar. Fable
5.1 entrega 9,5 puntos más que la referencia por un 8 por ciento más de coste;
Luna cuesta un 95 por ciento menos y entrega 3,5 puntos menos. Entre esos dos
extremos hay un factor de veinte en precio y trece puntos de índice. No existe
la pregunta \enquote{cuál es mejor} sin especificar antes cuánto vale la tarea
para el cliente.

\textbf{Segunda: los dos únicos puntos dominados son los penúltimos pasos de cada
compañía.} Claude Fable 5 está dominado por Claude Fable 5.1, que sale más barato
y puntúa más; GPT-5.6 Sol al precio de lanzamiento está dominado por su propio
precio promocional. Esto no es un defecto del gráfico, es el retrato de un mercado
donde el producto de hace ocho semanas ya no es comprable. Y sostiene una
recomendación operativa concreta: en este mercado no se firman contratos anuales
a precio cerrado.

\textbf{Tercera: en el perfil de agente, Claude Opus 5 domina a GPT-5.6 Sol.}
Dos puntos más de índice y un 29 por ciento menos de coste por tarea. Es la
afirmación más fuerte de todo el trabajo y por eso conviene acotarla de
inmediato: depende por completo del perfil elegido. Con el perfil de
conversación de la Figura~\ref{fig:precios}, GPT-5.6 Sol cuesta un 20 por ciento
menos que Claude Opus 5 y la dominancia desaparece. \textbf{Quien elija el perfil,
elige el ganador}, y por eso he publicado el perfil antes que el resultado.

\textbf{Cuarta: los vectores de Anthropic son verticales y los de OpenAI son
horizontales.} A1 y A2 suben cinco y cuatro puntos; O2, O3 y O4 recorren
34, 45 y 47 puntos porcentuales de coste sin mover el valor o bajándolo. Es la
firma geométrica de las dos estrategias del Cuadro~\ref{tab:clasificacion}: la
ofensiva de profundidad compra capacidad, la ofensiva de amplitud compra
cobertura de mercado. La única excepción es reveladora: A1 y O1 son casi el
mismo vector, porque las dos compañías sacaron su generación nueva con un mes de
diferencia y al precio de la anterior. Cuando compiten en lo mismo, se mueven
igual; la divergencia empieza en el paso siguiente.

\textbf{Quinta: el vector A2 es el más eficiente de los siete.} Sube cuatro
puntos de índice y baja veintinueve puntos porcentuales de coste a la vez, y no
procede de un modelo nuevo sino de un cambio de tarifa en la lectura de caché.
La lección para un CTO es incómoda y vale para cualquier sector: una parte
sustancial de la mejora de la relación entre valor y coste que percibe el cliente
no viene de la tecnología del producto, sino del diseño de la tarifa.

\section{Cómo se reparten el mercado}
\label{sec:mercado}

\begin{figure}[tbp]
  \centering
  \begin{tikzpicture}
    % ---------------- panel A: cuota por segmento ----------------
    \node[font=\scriptsize\itshape, gray!35!black, right] at (-4.35,3.10)
      {Panel A. Cuota por segmento, en por ciento del gasto o de la audiencia};
    \foreach \xv/\lab in {0/0, 2.4/20, 4.8/40, 7.2/60, 9.6/80, 12.0/100} {
      \draw[gray!22] (\xv,-0.35) -- (\xv,2.70);
      \node[below, font=\tiny, gray!55!black] at (\xv,-0.35) {\lab};
    }
    \foreach \y/\lab/\a/\b/\c in {%
      2.20/{Gasto empresarial en API\\de modelos (dic. 2025)}/4.80/3.24/2.52,
      1.10/{Gasto empresarial en\\codificación (dic. 2025)}/6.48/2.52/0.00,
      0.00/{Audiencia de asistentes\\de consumo (may. 2026)}/1.236/5.568/3.324} {
      \node[left, text width=4.2cm, align=right, font=\scriptsize]
        at (-0.15,\y) {\lab};
      \pgfmathsetmacro{\ab}{\a+\b}
      \pgfmathsetmacro{\abc}{\a+\b+\c}
      \fill[cAnt] (0,\y-0.24) rectangle (\a,\y+0.24);
      \fill[cOAI] (\a,\y-0.24) rectangle (\ab,\y+0.24);
      \fill[cGoo] (\ab,\y-0.24) rectangle (\abc,\y+0.24);
      \fill[gray!35] (\abc,\y-0.24) rectangle (12.0,\y+0.24);
    }
    \node[font=\tiny, white] at (2.40,2.20) {Anthropic 40};
    \node[font=\tiny, white] at (6.42,2.20) {OpenAI 27};
    \node[font=\tiny, white] at (9.30,2.20) {Google 21};
    \node[font=\tiny, white] at (3.24,1.10) {Anthropic 54};
    \node[font=\tiny, white] at (7.74,1.10) {OpenAI 21};
    \node[font=\tiny, gray!25!black] at (10.50,1.10) {resto 25};
    \node[font=\tiny, white] at (4.02,0.00) {OpenAI 46,4};
    \node[font=\tiny, white] at (8.47,0.00) {Google 27,7};
    \node[font=\tiny, white] at (0.62,0.00) {10,3};
    \node[font=\tiny, gray!25!black] at (11.28,2.20) {resto 12};
    \node[font=\tiny, gray!25!black] at (11.06,0.00) {resto 15,6};

    % ---------------- panel B: monetización ----------------
    \begin{scope}[shift={(0,-5.05)}]
      \node[font=\scriptsize\itshape, gray!35!black, right] at (-4.35,2.85)
        {Panel B. Tres razones de monetización, cada fila con su propia escala};
      \foreach \y/\lab/\xa/\xo/\va/\vo in {%
        2.00/{Ingreso anualizado,\\miles de millones de \$}/9.00/5.54/{65}/{40},
        1.00/{Ingreso por empleado,\\millones de \$}/9.00/3.08/{26,0}/{8,9},
        0.00/{Ingreso por punto de cuota de\\consumo, miles de millones de
          \$}/9.00/1.23/{6,31}/{0,86}} {
        \node[left, text width=4.2cm, align=right, font=\scriptsize]
          at (-0.15,\y) {\lab};
        \fill[cAnt] (0,\y+0.04) rectangle (\xa,\y+0.26);
        \fill[cOAI] (0,\y-0.26) rectangle (\xo,\y-0.04);
        \node[right, font=\tiny, cAnt] at (\xa+0.1,\y+0.15) {\va};
        \node[right, font=\tiny, cOAI] at (\xo+0.1,\y-0.15) {\vo};
      }
      \fill[cAnt] (0.15,-0.95) rectangle (0.70,-0.75);
      \node[right, font=\tiny, cAnt] at (0.78,-0.85) {Anthropic};
      \fill[cOAI] (2.60,-0.95) rectangle (3.15,-0.75);
      \node[right, font=\tiny, cOAI] at (3.23,-0.85) {OpenAI};
      \node[right, text width=6.6cm, font=\tiny, align=left, gray!25!black]
        at (5.05,-0.85) {la tercera fila divide el ingreso anualizado entre la
        cuota de audiencia de consumo del Panel A: mide cuánto dinero extrae
        cada compañía por unidad de atención};
    \end{scope}
  \end{tikzpicture}
  \caption{El mercado no se reparte, se divide en dos. Panel A: cuotas del gasto
  empresarial en interfaces de programación de modelos y del subsegmento de
  codificación según la encuesta de \num{495} responsables de decisión de
  noviembre de 2025 \parencite{menlo2025}, y cuota de audiencia de asistentes de
  consumo a finales de mayo de 2026 \parencite{sensortower2026}. Panel B: el
  ingreso anualizado de Anthropic en julio de 2026
  \parencite{techcrunch2026ingresos} y el de OpenAI en el mismo mes
  \parencite{bloomberg2026openai}, divididos por las plantillas declaradas
  \parencite{wikipedia2026anthropic,wikipedia2026openai} y por la cuota de
  audiencia del Panel A. Las dos últimas razones son cálculo propio.
  Elaboración propia.}
  \label{fig:mercado}
\end{figure}

La respuesta a la pregunta del enunciado es que \textbf{no se reparten un
mercado: cada uno domina el suyo y los dos están creciendo}. El Panel A de la
Figura~\ref{fig:mercado} lo muestra en tres franjas.

\textbf{En el gasto empresarial por interfaz de programación manda Anthropic.}
Cuarenta por ciento frente a veintisiete de OpenAI y veintiuno de Google, sobre
un gasto de \num{12500} millones de dólares en 2025. La serie histórica es más
elocuente que la foto: OpenAI tenía el 50 por ciento de ese gasto a finales de
2023 y Anthropic el 12 \parencite{menlo2025}. La inversión se produjo en dos
años.

\textbf{En codificación la asimetría es mayor.} Cincuenta y cuatro por ciento
frente a veintiuno \parencite{menlo2025}. Es el segmento donde Anthropic ha
concentrado su ofensiva de profundidad desde 2025, y es también el que explica
la diferencia de ingresos: la codificación empresarial es cara, recurrente y
consume sesiones largas, que es exactamente el perfil donde su estructura de
precios gana.

\textbf{En el consumo manda OpenAI, y no por poco.} Cuarenta y seis con cuatro
por ciento de la audiencia de asistentes a finales de mayo de 2026, frente al
27,7 de Google Gemini y al 10,3 de Claude \parencite{sensortower2026}. Conviene
matizar dos cosas. La primera es que ese 46,4 supone haber bajado del 50 por
primera vez en marzo de 2026 \parencite{techcrunch2026cuota}. La segunda es que
el segundo lugar no es de Anthropic sino de Google, lo que recuerda que este
mercado no es un duopolio aunque el foro lo plantee como tal.

El Panel B añade la lectura económica de ese reparto y produce el contraste más
llamativo del trabajo. Anthropic factura 65.000 millones anualizados con el 10,3
por ciento de la audiencia de consumo; OpenAI factura 40.000 millones con el
46,4 por ciento. Por cada punto de cuota de atención, Anthropic extrae 6.310
millones de dólares y OpenAI 860, una razón de 7,3 a 1. Son dos negocios
distintos vendiendo la misma tecnología: uno cobra por trabajo terminado y el
otro por atención, y el segundo ha empezado a monetizar esa atención con
publicidad, que en agosto de 2026 rondaba los mil millones anualizados
\parencite{tnw2026friar}.

Y sin embargo, la frontera se está borrando por el lado que más importa. El 14 de
agosto de 2026 la directora financiera de OpenAI comunicó a los accionistas que
los ingresos de empresa habían superado por primera vez a los de consumo, dos
trimestres antes de lo que ella misma había previsto, partiendo de un reparto de
60 a 40 a comienzos de año \parencite{tnw2026friar}. Es decir: \textbf{OpenAI está
entrando en el mercado donde Anthropic es fuerte, mientras Anthropic no está
entrando en el mercado donde OpenAI lo es}. Esa asimetría de movimientos es, a mi
juicio, el hecho competitivo más relevante de 2026 y el que hay que vigilar en
2027.

\begin{table}[tbp]
  \centering
  \small
  \setlength{\tabcolsep}{4.5pt}
  \caption{Ficha comparativa a 1 de septiembre de 2026.}
  \label{tab:ficha}
  \begin{tabularx}{\textwidth}{@{}>{\raggedright\arraybackslash}p{0.20\textwidth}
    >{\raggedright\arraybackslash}X >{\raggedright\arraybackslash}X@{}}
    \toprule
    & \textbf{Anthropic} & \textbf{OpenAI} \\
    \midrule
    Fundación & enero de 2021, ocho procedentes de OpenAI
      & diciembre de 2015, como entidad sin ánimo de lucro \\
    Forma jurídica & sociedad de beneficio público con fideicomiso de
      beneficio a largo plazo
      & sociedad de beneficio público desde 2025; la fundación conserva el
        26 por ciento \\
    Plantilla (2026) & del orden de \num{2500} & del orden de \num{4500} \\
    Última valoración & \num{965000} millones de dólares (mayo de 2026)
      & \num{852000} millones de dólares (abril de 2026) \\
    Ingreso anualizado & \num{65000} millones (julio de 2026)
      & \num{40000} millones (julio de 2026) \\
    Salida a bolsa & solicitud confidencial el 1 de junio de 2026
      & solicitud el 8 de junio de 2026 \\
    Modelo insignia & Claude Fable 5.1 y Mythos 5.1 (1 de septiembre)
      & GPT-5.6 Sol (9 de julio) \\
    Tarifa insignia & 10 y 50 dólares por millón de tokens; caché al 2,5\,\%
      & 4 y 20 dólares en contexto corto; 8 y 30 en contexto largo \\
    Ventana de contexto & 1 millón de tokens a tarifa plana
      & tramo de tarifa por encima de \num{272000} tokens \\
    Canal propio & ninguno: distribuye por AWS, Google Cloud y Microsoft
      Foundry
      & ChatGPT, aplicación de escritorio, publicidad, comercio y dispositivo
        anunciado \\
    \bottomrule
  \end{tabularx}
\end{table}

\section{Precisiones y limitaciones}

\begin{enumerate}
  \item \textbf{La elección del perfil de tarea determina el resultado del
    gráfico vectorial.} Con el perfil de agente, Claude Opus 5 domina a GPT-5.6
    Sol; con el perfil de conversación, no. He publicado los dos perfiles y los
    dos resultados precisamente para que el lector pueda desmontar mi
    conclusión. Ninguno de los dos perfiles procede de una medición de tráfico
    real: son construcciones mías, razonables pero arbitrarias.
  \item \textbf{Los costes están calculados en tokens, no en trabajo.} El
    tokenizador de Claude 4.7 y posteriores produce alrededor de un 30 por
    ciento más de tokens para el mismo texto que las generaciones anteriores de
    la propia Anthropic \parencite{anthropic2026precios}. No dispongo de una
    medición independiente que compare el tokenizador de Claude con el de GPT
    sobre el mismo corpus, de modo que la comparación de coste entre las dos
    compañías arrastra un error de magnitud desconocida.
  \item \textbf{Las versiones del índice compuesto no son comparables entre
    sí.} Claude Opus 4.8 figuraba con 61,4 puntos en mayo de 2026
    \parencite{aa2026opus48} y con 57 en la versión vigente
    \parencite{aa2026comparacion}. Todos los valores del eje vertical proceden de
    la misma versión, pero eso significa que el gráfico caduca con la siguiente.
  \item \textbf{Las puntuaciones de Terra y Luna están ajustadas por mí.} Sumo
    dos puntos a los valores publicados el 9 de julio \parencite{aa2026gpt56}
    para llevarlos a la versión vigente del índice, que es el desplazamiento
    observado en los dos modelos de los que sí tengo la doble medición. Es una
    interpolación, no un dato.
  \item \textbf{Los datos de cuota empresarial son de diciembre de 2025 y los de
    consumo de mayo de 2026.} La encuesta de Menlo Ventures tiene \num{495}
    respuestas de responsables de decisión estadounidenses \parencite{menlo2025};
    la de Sensor Tower mide audiencia de aplicaciones y web
    \parencite{sensortower2026}. Ni una ni otra es una medición de facturación
    auditada, y las dos tienen entre nueve y tres meses de retraso respecto de
    todo lo demás que aparece en este documento.
  \item \textbf{Las cifras de ingresos no son comparables sin reservas.} Las dos
    compañías son privadas, publican ritmos anualizados y no estados financieros
    auditados, y la propia fuente que da los 65.000 millones advierte de que las
    dos pueden calcular la métrica de forma distinta
    \parencite{techcrunch2026ingresos}. Las razones de monetización del Panel B
    de la Figura~\ref{fig:mercado} heredan toda esa incertidumbre y además la de
    los recuentos de plantilla, que proceden de fuentes terciarias.
  \item \textbf{Los datos de retención de talento son de 2025.} La edición de
    2026 del mismo informe ya no los desglosa por compañía
    \parencite{signalfire2025,signalfire2026}, de modo que marcan una dirección
    y no una situación actual.
  \item \textbf{He excluido a Google del análisis y no debería pasar
    desapercibido.} Con el 21 por ciento del gasto empresarial y el 27,7 de la
    audiencia de consumo, Gemini es el segundo jugador del mercado de consumo y
    el tercero del empresarial \parencite{menlo2025,sensortower2026}. El foro
    pedía dos competidores; el mercado tiene tres.
  \item \textbf{Varios hitos de producto proceden de fuentes terciarias.} Las
    cronologías de la Figura~\ref{fig:cronologia} se apoyan en recopilaciones
    enciclopédicas \parencite{wikipedia2026claude,wikipedia2026gpt56}, que he
    contrastado con los anuncios oficiales cuando estaban disponibles pero que
    tienen menos solidez que el resto de las fuentes empleadas.
\end{enumerate}

\section*{Conclusiones}

\textbf{El sector.} Modelos fundacionales de lenguaje vendidos como servicio, un
mercado donde la oferta, el precio y la propia definición del producto dependen
del estado de desarrollo del sistema tecnológico de dos compañías que, además,
comparten origen.

\textbf{La clasificación.} Las dos son ofensivas en el sentido de
\textcite{freeman1997}, que es la posición más cara de sostener y la que menos
margen deja para la imitación. Lo que las separa es la dirección de la ofensiva:
Anthropic va en profundidad sobre el trabajo agente y regala el estándar que le
da distribución sin poseer el canal; OpenAI va en amplitud sobre todos los
públicos y compra el canal entero, del navegador al dispositivo. En términos de
\textcite{teece1986}, una apuesta por el activo tecnológico y la otra por los
activos complementarios.

\textbf{El gráfico vectorial.} La frontera de eficiencia tiene cuatro puntos y
está repartida dos a dos entre las dos compañías. Los vectores de Anthropic son
verticales, ganan valor; los de OpenAI son horizontales, ganan coste. Los dos
únicos puntos dominados son los penúltimos pasos de cada una, lo que dice mucho
sobre la vida útil de cualquier análisis en este mercado. Y el vector más
eficiente de los siete no procede de un modelo nuevo, sino de un cambio de tarifa
en la lectura de caché.

\textbf{El reparto.} No se reparten un mercado, se reparten dos. Anthropic tiene
el 40 por ciento del gasto empresarial y el 54 de la codificación; OpenAI tiene
el 46,4 por ciento de la audiencia de consumo y ha convertido esa audiencia en
publicidad y comercio. Por cada punto de cuota de atención, Anthropic extrae 7,3
veces más ingreso. Pero el movimiento no es simétrico: en agosto de 2026 OpenAI
cruzó por primera vez el punto en el que su negocio de empresa supera al de
consumo, mientras Anthropic sigue sin entrar en el consumo. Quien quiera predecir
2027 tiene que mirar ahí.

La lección que me llevo como CTO no es sobre modelos de lenguaje.
\textbf{Cuando dos competidores comparten tecnología, origen y ritmo de
innovación, lo que decide el reparto del mercado no es la capacidad del producto
sino la arquitectura del precio y la elección de qué activo complementario se
posee.} Anthropic y OpenAI lanzan al mismo ritmo, puntúan casi igual en las
evaluaciones públicas y salen del mismo edificio. Se reparten el mercado como se
lo reparten porque uno cobra la ventana de contexto plana y regala el protocolo,
y el otro cobra el contexto por tramos y se queda con el canal. Esas dos
decisiones no las tomó ningún investigador.

\printbibliography[title={Referencias bibliográficas}]

\end{document}
